% ============ 27 Weil 猜想之证明(Riemann 猜想除外) (原书 pp.154-159) ============
\chapter{Weil 猜想之证明(Riemann 猜想除外)}\label{sec:27}
\index{Frobenius 映射}
\index{$\zeta$ 函数}
\index{有理函数}
\index{函数方程}
\index{Betti 数}

\subsection*{Frobenius 映射}

设 $A_0$ 是仿射 $\mathbb{F}_q$-代数.则 $a \mapsto a^{q}$ 是 $\mathbb{F}_q$-代数的同态 $f_0 \colon A_0 \to A_0$.通过标量扩张，得到 $\mathbb{F}$-代数的同态 $f \colon A \to A$，其中 $A = A_0 \otimes_{\mathbb{F}_q} \mathbb{F}$.相应的正则映射 $F \colon \operatorname{Specm} A \to \operatorname{Specm} A$ 称为 Frobenius 映射.

对 $\mathbb{F}_q$ 上的簇 $X_0$，Frobenius 映射 $F \colon X \to X$ 是唯一的正则映射，使得对每个开仿射 $U_0 \subset X_0$，$F(U) \subset U$ 且 $F|_{U}$ 是 $U$ 上的 Frobenius 映射.

容易检验：
\begin{enumerate}
  \item Frobenius 映射 $F \colon \mathbb{A}^{n} \to \mathbb{A}^{n}$ 是 $(t_1, \ldots, t_n) \mapsto (t^{q}_1, \ldots, t^{q}_n)$；
  \item Frobenius 映射 $F \colon \mathbb{P}^{n} \to \mathbb{P}^{n}$ 是 $(t_0 \colon \ldots \colon t_n) \mapsto (t^{q}_0 \colon \ldots \colon t^{q}_n)$；
  \item 对 $\mathbb{F}_q$ 上的任何正则映射 $\varphi' \colon Y_0 \to X_0$，图
  \[
  \begin{tikzcd}
  Y \arrow[r, "\varphi"] \arrow[d, "F"'] & X \arrow[d, "F"] \\
  Y \arrow[r, "\varphi"'] & X
  \end{tikzcd}
  \]
  交换.
\end{enumerate}
把这些论断结合起来，看到 Frobenius 映射在 $\mathbb{A}^{n}$ 的任何子簇上按 $(t_1, \ldots, t_n) \mapsto (t^{q}_1, \ldots, t^{q}_n)$ 作用，在 $\mathbb{P}^{n}$ 的任何子簇上按 $(t_0 \colon \ldots \colon t_n) \mapsto (t^{q}_0 \colon \ldots \colon t^{q}_n)$ 作用.

\begin{lemma}\label{lem:27.1}
Frobenius 映射 $F \colon X \to X$ 的次数为 $q^{\dim X}$(事实上，它是这个次数的有限映射).
\end{lemma}

\begin{proof}
在 $\mathbb{A}^{n}$ 上，$F$ 对应于 $\mathbb{F}$-代数的同态
\[
T_i \mapsto T^{q}_i \colon \mathbb{F}[T_1, \ldots, T_n] \to \mathbb{F}[T_1, \ldots, T_n].
\]
该同态的像是 $\mathbb{F}[T^{q}_1, \ldots, T^{q}_n]$，而 $\mathbb{F}[T_1, \ldots, T_n]$ 在 $\mathbb{F}[T^{q}_1, \ldots, T^{q}_n]$ 上自由，秩为 $q^{n}$(以元素 $T^{i_1}_1 \cdots T^{i_n}_n$，$0 \leq i_j \leq q - 1$ 为基).类似地，$\mathbb{F}(T_1, \ldots, T_n)$ 在其子域 $\mathbb{F}(T^{q}_1, \ldots, T^{q}_n)$ 上的次数为 $q^{n}$.这表明 $F \colon \mathbb{A}^{n} \to \mathbb{A}^{n}$ 是次数 $q^{n}$ 的有限映射.

在一般情形中，选取 $X_0$ 的函数域 $\mathbb{F}_q(X_0)$ 的超越基 $T_1, \ldots, T_n$.设 $f \colon \mathbb{F}(X) \to \mathbb{F}(X)$ 是由 $F \colon X \to X$ 定义的同态.则 $f(\mathbb{F}(X) \cap \mathbb{F}(T_1, \ldots)) = f(\mathbb{F}(T_1, \ldots))$ 且 $\mathbb{F}(X) = f(\mathbb{F}(X)) \cdot \mathbb{F}(T_1, \ldots)$，从而
\[
[\mathbb{F}(X) \colon f(\mathbb{F}(X))] = [\mathbb{F}(T_1, \ldots) \colon f(\mathbb{F}(T_1, \ldots))] = q^{n},
\]
这表明 $F$ 的次数为 $q^{n}$.
\end{proof}

\begin{lemma}\label{lem:27.2}
$F$ 在 $X$ 上的不动点就是 $X_0$ 中坐标在 $\mathbb{F}_q$ 中的点.每个不动点在 $(\Gamma_{F} \cdot \Delta)$ 中重数为 1.
\end{lemma}

\begin{proof}
$\mathbb{F}$ 的元素 $a$ 在 $\mathbb{F}_q$ 中，当且仅当 $a^{q} = a$，因此由 $F$ 按点的坐标的描述显然看到 $X^{F} = X(\mathbb{F}_q)$.

我断言 $F$ 的任何不动点 $P$ 处 $(dF)_P = 0$，从而 $F$ 满足引理~\ref{lem:25.6} 的条件.证明中可以把 $X_0$ 换成 $P$ 的仿射邻域 $U_0$，设 $U_0 = \operatorname{Specm} A_0$，$A_0 = \mathbb{F}_q[t_1, \ldots, t_n] = \mathbb{F}_q[T_1, \ldots, T_n]/\mathfrak{a}$.则 $t_i \circ F = t^{q}_i$，从而 $(dt_i)_P \circ (dF)_P = (dt^{q}_i)_P = qt^{q-1}_i (dt_i)_P = 0$，正如断言.
\end{proof}

\subsection*{\texorpdfstring{$N_m$ 表示为迹}{N_m 表示为迹}}

\begin{proposition}\label{prop:27.3}
设 $X_0$ 是 $\mathbb{F}_q$ 上的完备非奇异簇.对任何 $m$，
\[
N_m = \sum_{r} (-1)^{r} \operatorname{Tr}(F^{m}|H^{r}(X, \mathbb{Q}_{\ell})).
\]
\end{proposition}

\begin{proof}
由引理~\ref{lem:27.2}，看到 $(\Gamma_{F} \cdot \Delta)$ 是 $X_0$ 中坐标在 $\mathbb{F}_q$ 中的点数 $N_1$，按 Lefschetz 不动点公式 (\ref{thm:25.1})，它等于 $\sum_{r} (-1)^{r} \operatorname{Tr}(F|H^{r}(X, \mathbb{Q}_{\ell}))$.为得到一般情形，注意 $F^{m}$ 是 $X$ 相对于 $X_{0; \mathbb{F}_{q^m}}$ 的 Frobenius 映射，从而 $X^{F^{m}} = X(\mathbb{F}_{q^m})$.
\end{proof}

\begin{remark}\label{rem:27.4}
(a) 设 $X_0$ 是 $\mathbb{F}_q$ 上的完备非奇异簇.则 $F \colon X \to X$ 定义映射 $F^{*} \colon H^{r}(X, \mathbb{Q}_{\ell}) \to H^{r}(X, \mathbb{Q}_{\ell})$，且(见 \S 24)映射 $F_{*} \colon H^{r}(X, \mathbb{Q}_{\ell}) \to H^{r}(X, \mathbb{Q}_{\ell})$.因为 $F$ 是次数 $q^{\dim X}$ 的有限映射，复合 $F_{*} \circ F^{*} = q^{d}$(24.2f).

(b) 设 $\varphi' \colon Y_0 \to X_0$ 是 $\mathbb{F}_q$ 上完备非奇异簇的正则映射.则
\[
F^{*}\varphi^{*} = q^{\dim X - \dim Y}\varphi^{*}F^{*}.
\]
为证明它，对两边应用 $F_{*}$.左边得到
\[
F_{*}F^{*}\varphi^{*} = q^{\dim X}\varphi^{*},
\]
右边得到
\[
q^{\dim X - \dim Y}F_{*}\varphi^{*}F^{*} = q^{\dim X - \dim Y}\varphi^{*}F_{*}F^{*} = q^{\dim X}\varphi^{*}.
\]
注意这个论证对 $\mathbb{Z}/\ell^{n}\mathbb{Z}$ 系数($\ell \neq p$)的上同调群也成立.
\end{remark}

\subsection*{有理函数}

需要一个初等引理.

\begin{lemma}\label{lem:27.5}
把域 $k$ 上向量空间 $V$ 的自同态 $\varphi \colon V \to V$ 的特征多项式定义为
\[
P_{\varphi}(t) = \det(1 - \varphi t|V).
\]
若 $P_{\varphi}(t) = \prod (1 - c_i t)$，则 $\operatorname{Tr}(\varphi^{m}|V) = \sum c^{m}_i$.因此
\[
\log \frac{1}{P_{\varphi}(t)} = \sum_{m=1}^{\infty} \operatorname{Tr}(\varphi^{m}|V) \frac{t^{m}}{m}
\]
($k[[t]]$ 中的等式).
\end{lemma}

\begin{proof}
在可能扩张 $k$ 后，可以假定存在一组基，使 $\varphi$ 的矩阵为上三角
\[
\begin{pmatrix}
c_1 & *\\
& \ddots\\
0 & & c_d
\end{pmatrix}.
\]
相对于这组基，$\varphi^{m}$ 的矩阵为
\[
\begin{pmatrix}
c^{m}_1 & *\\
& \ddots\\
0 & & c^{m}_d
\end{pmatrix},
\]
论断由它显然.

对
\[
\log \frac{1}{1 - c_i t} = \sum_{m=1}^{\infty} c^{m}_i \frac{t^{m}}{m}
\]
两边关于 $i$ 求和，得到第二个公式.
\end{proof}

\begin{theorem}\label{thm:27.6}
对 $\mathbb{F}_q$ 上任何完备非奇异簇 $X_0$ 的维数 $d$；
\[
Z(X_0, t) = \frac{P_1(X_0, t) \cdots P_{2d-1}(X_0, t)}{P_0(X_0, t) \cdots P_{2d}(X_0, t)},
\]
其中
\[
P_r(X, t) = \det(1 - Ft|H^{r}(X, \mathbb{Q}_{\ell})).
\]
\end{theorem}

\begin{proof}
我们有
\[
Z(X_0, t) = \exp\left( \sum_{m} N_m \frac{t^{m}}{m} \right) \quad \text{(定义)}
\]
\[
= \exp\left( \sum_{m} \left( \sum_{r=0}^{2d} (-1)^{r} \operatorname{Tr}(F^{m}|H^{r}) \right) \frac{t^{m}}{m} \right) \quad \text{(\ref{prop:27.3})}
\]
\[
= \prod_{r=0}^{2d} \left( \exp\left( \sum_{m} \operatorname{Tr}(F^{m}|H^{r}) \frac{t^{m}}{m} \right) \right)^{(-1)^{r}} \quad \text{(把内层和移到外面)}
\]
\[
= \prod_{r=0}^{2d} P_r(t)^{(-1)^{r+1}} \quad \text{(\ref{lem:27.5}).}
\]
\end{proof}

\begin{remark}\label{rem:27.7}
因为 $F$ 在 $H^{0}(X, \mathbb{Q}_{\ell})$ 上按 1 作用，$P_0(X_0, t) = 1 - t$；因为 $F$ 在 $H^{2d}(X, \mathbb{Q}_{\ell})$ 上按 $q^{d}$ 作用，$P_{2d}(X_0, t) = 1 - q^{d}t$.一般地，$P_r(X_0, t) = 1 + \cdots \in \mathbb{Q}_{\ell}[t]$.
\end{remark}

\begin{corollary}\label{cor:27.8}
幂级数 $Z(X_0, t)$ 是系数在 $\mathbb{Q}$ 中的有理函数，即它属于 $\mathbb{Q}(t)$.
\end{corollary}

\begin{proof}
注意
\[
Z(X_0, t) \overset{\mathrm{def}}{=} \exp\left( \sum_{m \geq 1} N_m t^{m}/m \right)
\]
是系数在 $\mathbb{Q}$ 中的幂级数.定理~\ref{thm:27.6} 表明它是系数在 $\mathbb{Q}_{\ell}$ 中的有理函数，按下一个引理，这蕴含它是系数在 $\mathbb{Q}$ 中的有理函数.
\end{proof}

\begin{lemma}\label{lem:27.9}
设 $k \subset K$ 是域，$f(t) \in k[[t]]$；若 $f(t) \in K(t)$，则 $f(t) \in k(t)$.
\end{lemma}

\begin{proof}
由 Bourbaki, \emph{Alg\`ebre}, IV.5，习题 3：
设 $f(t) = \sum_{n=1}^{\infty} a_i t^{i}$ 是域 $k$ 上的形式幂级数.
\begin{enumerate}
  \item 幂级数 $f(t)$ 在 $k(t)$ 中，当且仅当存在 $k$ 中不全为零的有限序列 $\lambda_1, \ldots, \lambda_r$ 与整数 $d$，使得对所有 $n \geq d$，
  \[
  \lambda_1 a_n + \lambda_2 a_{n+1} + \cdots + \lambda_r a_{n+r-1} = 0.
  \]
  \item 设
  \[
  H^{(k)}_{n} = \begin{vmatrix}
  a_n & a_{n+1} & \cdots & a_{n+k-1}\\
  a_{n+1} & a_{n+2} & \cdots & a_{n+k}\\
  \vdots & \vdots & & \vdots\\
  a_{n+k-1} & a_{n+k} & \cdots & a_{n+2k-2}
  \end{vmatrix}
  \]
  (Hankel 行列式).若对某个 $d$ 与所有 $j \geq 0$，$H^{(r+1)}_{d+j} = 0$ 且 $H^{(r)}_{d+j} \neq 0$，则 $f(t) \in k(t)$.
  \item 证明
  \[
  H^{(k)}_{n} H^{(k)}_{n+2} - H^{(k+1)}_{n} H^{(k-1)}_{n+2}
  \]
  是 $H^{(k)}_{n+1}$ 的幂.由此推出若 $H^{(k+1)}_{m+j} = 0$($0 \leq j \leq s - 1$)，则 $s$ 个行列式 $H^{(k)}_{m+j}$($1 \leq j \leq s$)全为零或全非零.
  \item 由 (b) 与 (c) 推出 $f(X) \in k(t)$，当且仅当存在两个整数 $d$ 与 $r$，使得对所有整数 $j \geq 1$，$H^{(r+1)}_{d+j} = 0$.
\end{enumerate}
显然，若 (d) 中的条件在 $K$ 中满足，则在 $k$ 中也满足.
\end{proof}

该推论不蕴含多项式 $P_r(X_0, t)$ 有有理系数(更不蕴含它们与 $\ell$ 无关).它说的是，一旦去掉任何公因子，(\ref{thm:27.6}) 中表达式的分子与分母将是系数在 $\mathbb{Q}$ 中的多项式，且与 $\ell$ 无关.

\subsection*{整性}

设 $x$ 是 $X_0$ 的闭点.剩余域 $\kappa(x) \overset{\mathrm{def}}{=} \mathcal{O}_{X_0, x}/\mathfrak{m}_x$ 是 $\mathbb{F}_q$ 的有限扩张；我们设 $\deg x = [\kappa(x) \colon \mathbb{F}_q]$.

$X_0$ 中坐标在 $\mathbb{F}_{q^m}$ 中的点是映射 $\operatorname{Specm} \mathbb{F}_{q^m} \to X_0$.给出像为 $x$ 的这种映射与给出 $\mathbb{F}_q$-同态 $\kappa(x) \to \mathbb{F}_{q^m}$ 是一回事.$x$ 对 $N_m$ 的贡献 $N_m(x)$ 是这样的同态的个数.由有限域的理论(FT 4.26)，看到
\[
N_m(x) = \begin{cases}
\deg x & \text{若 } \deg x \mid m;\\
0 & \text{其他情形}.
\end{cases}
\]
回忆 $\log\left( \frac{1}{1-s} \right) = \sum_{m \geq 1} \frac{s^{n}}{n}$，因此
\[
\log \frac{1}{1 - t^{\deg x}} = \sum_{n \geq 1} \frac{t^{n \cdot \deg x}}{n}.
\]
该和中 $t^{m}/m$ 的系数，除非 $\deg x \mid m$(此时它为 $\deg x$)，否则为零.因此
\[
\log \frac{1}{1 - t^{\deg x}} = \sum_{m} N_m(x) \frac{t^{m}}{m}.
\]
对 $X_0$ 的所有闭点求和并取指数，发现
\begin{equation}
Z(X_0, t) = \prod_{x \in X_0} \frac{1}{1 - t^{\deg x}}. \tag{4}
\end{equation}
因此
\[
Z(X_0, t) \in 1 + t\mathbb{Z}[[t]].
\]
下一个引理中，$\ell$ 是任何素数(例如 $p$).

\begin{lemma}\label{lem:27.10}
设 $f(t) = g(t)/h(t)$，其中
\[
f(t) \in 1 + t \cdot \mathbb{Z}_{\ell}[[t]], \qquad g(t), h(t) \in 1 + t \cdot \mathbb{Q}_{\ell}[t].
\]
若 $g$ 与 $h$ 互素，则它们的系数在 $\mathbb{Z}_{\ell}$ 中.
\end{lemma}

\begin{proof}
必须证明 $g$ 与 $h$ 的系数有 $\ell$ 进绝对值 $\leq 1$.在可能把 $\mathbb{Q}_{\ell}$ 换成有限扩张域后，可以假定 $h(t)$ 分裂，设 $h(t) = \prod (1 - c_i t)$.若 $|c_i|_{\ell} > 1$，则 $|c^{-1}_i|_{\ell} < 1$，幂级数 $f(c^{-1}_i)$ 收敛.但
\[
f(t) \cdot h(t) = g(t) \Rightarrow f(c^{-1}_i) \cdot h(c^{-1}_i) = g(c^{-1}_i).
\]
因为 $h(c^{-1}_i) = 0$ 而 $g(c^{-1}_i) \neq 0$，这不可能.因此 $|c_i|_{\ell} < 1$.由于这对所有 $i$ 成立，$h(t) \in \mathbb{Z}_{\ell}[t]$.

因为 $f(t)^{-1} \in 1 + t \cdot \mathbb{Z}_{\ell}[[t]]$，把同样的论证应用于 $f(t)^{-1}$，表明 $g(t) \in \mathbb{Z}_{\ell}[t]$.
\end{proof}

\begin{proposition}\label{prop:27.11}
设
\[
Z(X_0, t) = \frac{P(t)}{Q(t)},
\]
其中 $P(t), Q(t) \in \mathbb{Q}[t]$ 互素($P$ 与 $Q$ 由 27.8 存在).当 $P(X_0, t)$ 与 $Q(X_0, t)$ 选取为常数项 1 时，它们的系数在 $\mathbb{Z}$ 中.
\end{proposition}

\begin{proof}
前一个引理的假设对所有素数 $\ell$(包括 $p$)都成立.因此，$P(X_0, t)$ 与 $Q(X_0, t)$ 的系数对所有 $\ell$ 都是 $\ell$ 进整数，这蕴含它们是整数.
\end{proof}

\subsection*{函数方程}

\begin{theorem}\label{thm:27.12}
对 $\mathbb{F}_q$ 上任何完备非奇异簇 $X_0$，
\[
Z\left( X_0, \frac{1}{q^{d}t} \right) = \pm q^{d\chi/2} t^{\chi} Z(X_0, t),
\]
其中 $\chi = \sum (-1)^{r}\beta_r = (\Delta \cdot \Delta)$.
\end{theorem}

\begin{proof}
考虑配对
\[
H^{2d-r}(X, \mathbb{Q}_{\ell}) \times H^{r}(X, \mathbb{Q}_{\ell}(d)) \to H^{2d}(X, \mathbb{Q}_{\ell}) \xrightarrow{\varphi_X} \mathbb{Q}_{\ell}, \qquad d = \dim X.
\]
由 $F_{*}$ 的定义，
\[
\varphi_X(F_{*}(x) \smile x') = \varphi_X(x \smile F_{*}(x')), \qquad x \in H^{2d-r}(X), \ x' \in H^{r}(X).
\]
因此，$F_{*}$ 在 $H^{r}(X)$ 上的特征值就是 $F_{*}$ 在 $H^{2d-r}(X)$ 上的特征值.但 $F_{*} = q^{d}/F^{*}$(见注记~\ref{rem:27.4}a)，因此若 $\alpha_1, \ldots, \alpha_s$ 是 $F^{*}$ 在 $H^{r}(X, \mathbb{Q}_{\ell})$ 上的特征值，则 $q^{d}/\alpha_1, \ldots, q^{d}/\alpha_s$ 是 $F^{*}$ 在 $H^{2d-r}(X, \mathbb{Q}_{\ell})$ 上的特征值.这蕴含论断.
\end{proof}

\begin{remark}\label{rem:27.13}
若 $d$ 为奇数，或 $q^{d/2}$ 作为 $F$ 在 $H^{d}(X, \mathbb{Q}_{\ell})$ 上的特征值出现偶数次，则符号为 $+$，否则为 $-$.
\end{remark}

\subsection*{小结}

\begin{quote}
\refstepcounter{theorem}\label{item:27.14}
\textbf{27.14} 设 $X_0$ 是 $\mathbb{F}_q$ 上的完备非奇异簇.

(a) 则 $Z(X_0, t) \in \mathbb{Q}(t)$，且满足函数方程
\[
Z\left( \frac{1}{q^{d}t} \right) = \pm q^{d\chi/2} \cdot t^{\chi} \cdot Z(t),
\]
其中 $\chi = (\Delta \cdot \Delta)$.

(b) 此外，对每个 $\ell \neq p$，有表达式
\[
Z(t) = \frac{P_{1, \ell}(t)P_{3, \ell}(t) \cdots P_{2d-1, \ell}(t)}{P_{0, \ell}(t)P_{2, \ell}(t) \cdots P_{2d, \ell}(t)},
\]
其中 $P_{0, \ell}(t) = 1 - t$，$P_{2d, \ell}(t) = 1 - q^{d}t$.

(c) 若对固定的 $\ell$，$P_{r, \ell}$ 两两互素，则对 $1 \leq r \leq 2d - 1$，
\[
P_{r, \ell}(t) = 1 + \sum a_{i, r}t^{i} \in \mathbb{Z}[t].
\]

(d) 若对所有 $\ell \neq p$，$P_{r, \ell}$ 的逆根都有绝对值 $q^{r/2}$，则 $P_{r, \ell}(t) \in 1 + t\mathbb{Z}[t]$ 且与 $\ell$ 无关.

(e) 设 $\beta_r = \deg P_r(t)$.则 $\chi = \sum_{r} (-1)^{r}\beta_r$，且若 $X_0$ 提升为特征零的簇 $X_1$，则 $\beta_r$ 是 $X_1$ 看作复数域上的簇时的 Betti 数.
\end{quote}

论断 (a) 在 (\ref{cor:27.8}),(\ref{thm:27.12}) 中证明，(b) 在 (\ref{thm:27.6}) 中证明.若 $P_{r, \ell}(t)$ 两两互素，则 (\ref{prop:27.11}) 表明
\[
\prod_{r \text{ 奇}} P_{r, \ell}(t) \in 1 + t\mathbb{Z}[t], \qquad \prod_{r \text{ 偶}} P_{r, \ell}(t) \in 1 + t\mathbb{Z}[t].
\]
因此 $P_{r, \ell}(t)$ 的逆根是代数整数，这蕴含 $P_{r, \ell}(t) \in 1 + t\mathbb{Z}[t]$，由此得到 (c).(d) 的假设蕴含 (c) 的假设，因此 $P_{r, \ell}(t)$ 有整数系数；此外，$P_{r, \ell}(t)$ 与 $\ell$ 无关地被刻画为 $Z(X_0, t)$ 的分子或分母中根有绝对值 $q^{r/2}$ 的因子.最后，(e) 由 25.1 与 20.5 推出.

下面的论断 27.15 蕴含 (c) 与 (d) 的假设成立，并完成 Weil 猜想的证明.

包含 $\mathbb{Q}$(例如 $\mathbb{Q}_{\ell}$)的某个域的元素 $\alpha$ 称为代数数，如果它是多项式 $P(T) \in \mathbb{Q}[T]$ 的根.可以选取 $P(T)$ 首一且不可约；则 $P(T)$ 在 $\mathbb{C}$ 中的根称为 $\alpha$ 的复共轭.等价地，$\alpha$ 代数，如果它生成 $\mathbb{Q}$ 的有限扩张 $\mathbb{Q}[\alpha]$；$\alpha$ 的复共轭是它在各种同态 $\mathbb{Q}[\alpha] \to \mathbb{C}$ 下的像.

\begin{theorem}\label{thm:27.15}
设 $X_0$ 是 $\mathbb{F}_q$ 上的非奇异射影簇.则 $F$ 在 $H^{r}(X, \mathbb{Q}_{\ell})$ 上的特征值是代数数，且所有复共轭都有绝对值 $q^{r/2}$.
\end{theorem}

代数数的所有复共轭都有相同绝对值是极不寻常的.例如，$1 + \sqrt{2}$ 没有这个性质.

课程的余下部分致力于 Deligne 对定理~\ref{thm:27.15} 的证明.

课程这一部分的讲义基于 Deligne 的原始文章：
\begin{quote}
Deligne, Pierre, \emph{La conjecture de Weil. I}, Inst. Hautes \'Etudes Sci. Publ. Math. No. 43 (1974), 273--307; MR 49 \#5013
\end{quote}
与
\begin{quote}
Katz, N., \emph{Deligne's Proof of the Weil-Riemann Conjectures}. Handwritten notes (by S. Bloch) of a course at Princeton University, 1973--74.
\end{quote}
